\documentclass[11pt,a4paper]{article}
\usepackage[italian]{babel}
\usepackage[margin=2.5cm]{geometry}
\usepackage{parskip}
\usepackage{amsmath, amssymb, amsthm}
\usepackage[most]{tcolorbox}
\usepackage{array}
\usepackage{tikz}
\usetikzlibrary{decorations.pathreplacing}
\usepackage{hyperref}

% --- Riquadri con bordo nero, senza numero ---
% Uso: \begin{definizione} ... \end{definizione}
%      \begin{teorema}[Nome] ... \end{teorema}  (il nome è facoltativo)
\tcbset{nota/.style={colback=white, colframe=black, boxrule=0.5pt,
  sharp corners}}
\NewTColorBox{definizione}{}{nota,
  before upper={\textbf{Definizione.}\ }}
\NewTColorBox{teorema}{o}{nota,
  before upper={\textbf{Teorema\IfValueT{#1}{ (#1)}.}\ }}
\NewTColorBox{proposizione}{o}{nota,
  before upper={\textbf{Proposizione\IfValueT{#1}{ (#1)}.}\ }}
\NewTColorBox{assioma}{o}{nota, boxrule=1pt,
  before upper={\textbf{Assioma\IfValueT{#1}{ (#1)}.}\ }}

% --- Ambienti semplici (senza riquadro) ---
\theoremstyle{definition}
\newtheorem*{esempio}{Esempio}
\newtheorem*{osservazione}{Osservazione}
\newtheorem*{esercizio}{Esercizio}

% V e F nelle tavole di verità
\newcommand{\V}{\textsf{V}}
\newcommand{\F}{\textsf{F}}

\title{Applicazioni e funzioni numeriche}
\author{Marco Cerbella}
\date{Lezione 3 -- 1 ottobre 2026}

\begin{document}
\maketitle

\section{Applicazioni tra insiemi}

\begin{definizione}
Siano $A$ e $B$ due insiemi. Si dice \emph{applicazione} da $A$ in $B$ una relazione che ad ogni elemento di $A$ associa \textbf{uno ed un solo} elemento di $B$. In simboli
\[
f : A \to B, \qquad a \mapsto f(a) = b.
\]
\end{definizione}

\begin{itemize}
    \item $f(a)$ è l'\emph{immagine} di $a$ tramite $f$;
    \item $A$ è il \emph{dominio} di $f$;
    \item $B$ è il \emph{codominio} di $f$;
    \item $f(A) = \{f(a) \mid a \in A\} \subseteq B$ è l'\emph{immagine di $f$}, indicata con $\mathrm{Im}(f)$;
    \item il \emph{grafico} di $f$ è
    $\mathrm{Gr}(f) = \{(a, b) \in A \times B : b = f(a),\ \forall\, a \in A\} = \{(a, f(a)) \mid a \in A\}$.
\end{itemize}

\begin{osservazione}
Da ogni elemento del dominio deve partire \textbf{una sola freccia}: non è ammesso che $a$ abbia due immagini diverse $b_1 \neq b_2$. Possono invece esistere elementi del codominio senza freccia in arrivo (non tutto $B$ è immagine), o con più frecce in arrivo.
\end{osservazione}

\begin{center}
\begin{tikzpicture}[every node/.style={font=\small}]
  \draw (0,0) ellipse (0.9 and 1.1);
  \draw (3.2,0) ellipse (1.2 and 1.1);
  \node at (0,1.35) {$A$}; \node at (3.2,1.35) {$B$};
  \fill (0,0.5) circle (1.5pt) node[left=2pt] {$a$};
  \fill (3.2,0.5) circle (1.5pt) node[right=2pt] {$f(a) = b$};
  \draw[->] (0,0.5) -- (3.2,0.5);
  \node at (1.6,0.75) {$f$};
  % non ammesso
  \begin{scope}[xshift=7.5cm]
    \draw (0,0) ellipse (0.9 and 1.1);
    \draw (3.2,0) ellipse (1.2 and 1.1);
    \node at (0,1.35) {$A$}; \node at (3.2,1.35) {$B$};
    \fill (0,-0.1) circle (1.5pt) node[left=2pt] {$a_1$};
    \fill (3.2,0.5) circle (1.5pt) node[right=2pt] {$b_1$};
    \fill (3.2,-0.7) circle (1.5pt) node[right=2pt] {$b_2$};
    \draw[->] (0,-0.1) -- (3.2,0.5);
    \draw[->] (0,-0.1) -- (3.2,-0.7);
    \draw[red, thick] (1.3,0.3) -- (1.9,-0.5);
    \draw[red, thick] (1.3,-0.5) -- (1.9,0.3);
    \node at (1.6,-1.5) {non è un'applicazione};
  \end{scope}
\end{tikzpicture}
\end{center}

\section{Funzioni numeriche}

\begin{definizione}
Si dice \emph{funzione} (numerica) una applicazione che ha come dominio un sottoinsieme $D$ di $\mathbb{R}$ e come codominio $\mathbb{R}$:
\[
f : D \subseteq \mathbb{R} \to \mathbb{R}, \qquad x \mapsto f(x).
\]
\end{definizione}

L'immagine è $f(D) = \{f(x) \mid x \in D\} \subseteq \mathbb{R}$ ($\mathrm{Im}(f)$) e il grafico è
\[
\mathrm{Gr}(f) = \{(x, y) \in \mathbb{R} \times \mathbb{R} : y = f(x),\ x \in D\}.
\]

La proprietà fondamentale, che ad ogni ingresso $x \in D$ corrisponde una ed una sola uscita $f(x) \in \mathbb{R}$, ha un significato geometrico: \textbf{ogni retta parallela all'asse delle ordinate che taglia l'asse delle ascisse in un punto $x \in D$ interseca il grafico di $f$ in uno ed un solo punto.} Una curva che viene tagliata da una retta verticale in due o più punti non è il grafico di una funzione.

\section{Dominio}

\begin{definizione}
Il \emph{dominio} di una funzione data da un'espressione è il più grande insieme delle variabili indipendenti per cui le espressioni coinvolte sono ben definite.
\end{definizione}

Le regole più frequenti per determinarlo:
\begin{enumerate}
    \item gli argomenti delle radici di ordine pari devono essere non negativi;
    \item gli argomenti dei logaritmi devono essere positivi;
    \item i denominatori devono essere non nulli.
\end{enumerate}

\begin{esempio}
$f(x) = \sqrt{x - 2}$ ha dominio $[2, +\infty)$; $f(x) = \ln(x + 1)$ ha dominio $(-1, +\infty)$; $f(x) = \dfrac{1}{x - 3}$ ha dominio $\mathbb{R} \setminus \{3\}$.
\end{esempio}

\end{document}
